\documentclass[11pt]{amsart}
\usepackage{geometry}                % See geometry.pdf to learn the layout options. There are lots.
\geometry{letterpaper}                   % ... or a4paper or a5paper or ... 
%\geometry{landscape}                % Activate for for rotated page geometry
%\usepackage[parfill]{parskip}    % Activate to begin paragraphs with an empty line rather than an indent
\usepackage{graphicx}
\usepackage{amssymb}
\usepackage{epstopdf}
\usepackage{color}
\usepackage{colortbl}
\DeclareGraphicsRule{.tif}{png}{.png}{`convert #1 `dirname #1`/`basename #1 .tif`.png}

\title{PS 2.7}
%\date{}                                           % Activate to display a given date or no date

\begin{document}
\maketitle
%\section{}
%\subsection{}

\noindent Haga un diagrama de flujo para calcular el precio del billete ida y vuelta en ferrocarril, conociendo la distancia del viaje de ida y el tiempo de estancia. Se sabe además que si el número de días de estancia es superior a 7 y la distancia total(ida y vuelta) a recorrer es superior a 800 km, el billete tiene una reducción del 30\%. El precio por km es de \$0.23.\\

\noindent \quad Datos: DIST, TIEM.\\

\noindent  Donde:\\

\noindent \quad \quad DIST es una variable de tipo entero que representa la distancia del viaje de ida.\\

\noindent \quad \quad TIEM es una variable de tipo entero que representa el tiempo de estancia.\\

\end{document}   